\section{The fruit and the door}

Augustine's former thief once enjoyed taking someone else's property. Now he enjoys giving what he owns to someone who lacks it. For a communicant who recognizes his own past in that account, the fruit of the Communion psalm includes the change in what he wants for his neighbor. Bellarmine identifies the fruit with Christ, whom Mary bore. A communicant who has been wondering how the servants can sustain their work has already received the Lord who gives them life.

At the Gospel the two questions have different urgency. Peter asleep after his brave assurances needs to learn how much he must ask of God if he is to do even what he wishes to do. The porter who has accepted a charge needs to find out whose safety depends on his attending to it. A congregation can contain both the person who scarcely dares to pray and the one who is carrying a considerable responsibility well. It can also contain someone who has allowed that very responsibility to feed his vanity. At the Collect all three can ask to be found doing right. What that would require of each will become apparent in the work awaiting them when they leave the church.
